\documentclass[12pt,a4paper,fontset=fandol]{ctexart}

% ==================== 页面与基础宏包 ====================
\usepackage[a4paper,margin=2.2cm]{geometry}
\usepackage{xcolor}
\usepackage{titlesec}
\usepackage{fancyhdr}
\usepackage{lastpage}
\usepackage{setspace}
\usepackage{enumitem}
\usepackage{amsmath,amssymb}
\usepackage{booktabs}
\usepackage{array}
\usepackage{longtable}
\usepackage{multirow}
\usepackage{float}
\usepackage{caption}
\usepackage{hyperref}
\usepackage[most]{tcolorbox}
\usepackage{abstract}

% ⭐ 关键：支持 svg
\usepackage{svg}
\svgsetup{inkscape=inkscape}

% ==================== 颜色定义 ====================
\definecolor{xjtuRed}{RGB}{165,30,55}
\definecolor{darkText}{RGB}{45,45,45}
\definecolor{lightGray}{RGB}{245,245,245}
\definecolor{midGray}{RGB}{120,120,120}
\definecolor{lineGray}{RGB}{220,220,220}

% ==================== 超链接 ====================
\hypersetup{
    colorlinks=true,
    linkcolor=xjtuRed,
    urlcolor=xjtuRed,
    citecolor=xjtuRed
}

% ==================== 行距 ====================
\setstretch{1.35}

% ==================== 页眉页脚 ====================
\pagestyle{fancy}
\fancyhf{}
\fancyhead[L]{\textcolor{midGray}{西安交通大学课程报告}}
\fancyhead[R]{\textcolor{midGray}{\leftmark}}
\fancyfoot[C]{\textcolor{midGray}{\thepage\ /\ \pageref{LastPage}}}

% ==================== 标题 ====================
\titleformat{\section}
  {\Large\bfseries\color{xjtuRed}}
  {\thesection}{0.8em}{}

\titleformat{\subsection}
  {\large\bfseries\color{darkText}}
  {\thesubsection}{0.8em}{}

% ==================== 信息 ====================
\newcommand{\ReportTitle}{课程报告题目}
\newcommand{\CourseName}{课程名称}
\newcommand{\StudentName}{你的名字}
\newcommand{\StudentID}{你的学号}
\newcommand{\School}{你的学院}
\newcommand{\Major}{你的专业}
\newcommand{\ClassName}{你的班级}
\newcommand{\Instructor}{你的指导教师}
\newcommand{\DateInfo}{\today}

% ==================== 封面 ====================
\newcommand{\makecover}{
\begin{titlepage}
\thispagestyle{empty}

\vspace*{0.5cm}

\noindent
    \begin{minipage}[c]{0.2\textwidth}
        \includegraphics[width=0.72\textwidth]{figs/logo.png}
    \end{minipage}
\hfill
\begin{minipage}[c]{0.75\textwidth}
    {\fontsize{22pt}{26pt}\selectfont\bfseries\color{xjtuRed} 西安交通大学}\\
    {\large\color{midGray} Xi'an Jiaotong University}
\end{minipage}

\vspace{2.2cm}

\begin{center}
    {\fontsize{26pt}{32pt}\selectfont\bfseries\color{darkText}\ReportTitle}

    \vspace{1.2em}

    \begin{tcolorbox}[
        colback=lightGray,
        colframe=lightGray,
        boxrule=0pt,
        arc=2mm,
        width=0.78\textwidth
    ]
    \centering
    {\Large\bfseries\color{xjtuRed} 课程报告}
    \end{tcolorbox}
\end{center}

\vspace{2cm}

\renewcommand{\arraystretch}{1.6}
\begin{center}
\begin{tabular}{p{3.2cm} p{8.5cm}}
课程名称： & \CourseName \\
学\hspace{1cm}院： & \School \\
专\hspace{1cm}业： & \Major \\
班\hspace{1cm}级： & \ClassName \\
姓\hspace{1cm}名： & \StudentName \\
学\hspace{1cm}号： & \StudentID \\
指导教师： & \Instructor \\
日\hspace{1cm}期： & \DateInfo \\
\end{tabular}
\end{center}

\vfill

\begin{center}
\color{midGray}
简约扁平风课程报告模板
\end{center}

\end{titlepage}
}

% ==================== 正文开始 ====================
\begin{document}

\makecover

\pagenumbering{Roman}

\begin{abstract}
本文围绕课程主题展开分析与总结，介绍研究背景、核心问题、方法设计、过程分析与结果讨论，并在最后给出个人理解与课程收获。该模板注重结构清晰、页面简洁和信息层次分明，适合用于本科阶段课程报告、实验报告、课程论文与读书报告的写作。
\end{abstract}

\vspace{1em}
\noindent\textbf{关键词：}课程报告；LaTeX；西安交通大学；模板设计

\tableofcontents
\newpage

\pagenumbering{arabic}

\section{引言}
这里写课程报告的研究背景、写作目的和主要内容。一般来说，这一部分主要回答三个问题：为什么要写这份报告，这份报告围绕什么问题展开，以及本文准备如何组织后续内容。写这一部分时，建议使用完整段落，不需要过度拆分成很多小点。

\section{理论基础}
这一部分用于介绍完成本报告所依赖的概念、原理和基本知识。对于计算机类课程，可以在这里介绍相关定义、系统结构、算法思想或者必要的数学基础。如果正文后面涉及符号和公式，也可以先在这一节统一说明。

\subsection{基本概念}
这里填写与你的主题相关的核心概念。例如：
\begin{equation}
T(n)=O(n\log n)
\end{equation}
这类公式可以根据课程内容自行替换。

\subsection{相关原理}
这里进一步说明问题背后的理论依据。若需要写多行公式，可以使用：
\begin{align}
A &= B + C \\
  &= D + E
\end{align}

\section{问题分析与方法设计}
这一部分是报告正文的核心区域。你可以在这里说明问题是什么、为什么这样分析、采用了什么思路、最后如何形成解决方案。若是课程论文，可以写你的理解与推导过程；若是实验报告，可以写实验方案和技术路线；若是课程设计，可以写系统模块与实现框架。

\subsection{整体思路}
这里介绍整体结构。例如，可以先说明目标，再解释为什么采用当前方案，最后说明各部分之间的联系。

\subsection{关键步骤}
这里写关键分析过程、实现步骤或设计逻辑。建议重点写“为什么这样做”，而不仅仅是“做了什么”。

\section{结果分析}
如果你的报告有实验、程序运行结果、案例分析或者仿真现象，可以统一放在这一节。重点不只是罗列现象，而是解释结果和课程知识之间的联系。

\subsection{结果展示}
下面是一个表格示例：
\begin{table}[H]
    \centering
    \caption{结果示例表}
    \begin{tabular}{cccc}
        \toprule
        编号 & 参数设置 & 结果A & 结果B \\
        \midrule
        1 & 配置一 & 0.91 & 0.84 \\
        2 & 配置二& 0.93 & 0.86 \\
        3 & 配置三 & 0.95 & 0.89 \\
        \bottomrule
    \end{tabular}
\end{table}

\subsection{分析讨论}
这里写你对结果的解释。例如，为什么某个配置更优、结果是否符合理论预期、哪些地方还有局限性，以及是否还有可以改进的空间。真正体现理解深度的，通常就是这一部分。

\section{总结与收获}
这里可以写你通过本次课程报告获得了哪些认识，课程内容如何帮助你理解某个问题，以及这次写作对后续学习的意义。若老师更偏向正式风格，可以写得客观一些；若允许加入个人思考，也可以适度写自己的理解。

\section*{参考文献}
\addcontentsline{toc}{section}{参考文献}

\begin{enumerate}[label={[}\arabic*{]}]
    \item 作者. 文献标题[J]. 期刊名, 年份, 卷(期): 页码.
    \item 作者. 书名[M]. 出版地: 出版社, 年份.
    \item 作者. 文献标题[EB/OL]. 链接.
\end{enumerate}

\end{document}